% ============================================================================
% Fichier  : partie4/P04_C18_systemes.tex
% Titre    : Sécurité des Systèmes
% Auteur   : MINKA MI NGUIDJOÏ Thierry Emmanuel
% Version  : 1.0 -- novembre 2026
% ============================================================================
% Sources : CIS Benchmarks et CIS Controls v8 (contrôles 4, 7, 8, 10) ;
% NIST SP 800-40 Rev. 4 (2022, gestion des correctifs) ; NIST SP 800-92
% (2006, gestion des journaux) ; NIST SP 800-123 (2008, sécurité des
% serveurs).
% ============================================================================

\chapter{Sécurité des Systèmes}
\label{chap:systemes}
\chaptermark{Sécurité des systèmes}

\epigraph{%
    « Un service inutile est une porte que l'on paie pour garder. »%
}{Principe du cours}

\niveauChapitre{Intermédiaire}{Chapitres~\ref{chap:architecture}
et~\ref{chap:iam}~; administration Windows et Linux.}

\begin{boxobjectifs}
\begin{itemize}
    \item Réduire la surface d'attaque d'un système par le durcissement.
    \item Appliquer des configurations de référence (CIS Benchmarks).
    \item Organiser la gestion des correctifs selon le risque.
    \item Définir ce qu'il faut journaliser, et comment.
    \item Situer l'antivirus, l'EDR et la liste blanche d'applications.
\end{itemize}
\end{boxobjectifs}

\bigskip

% ============================================================================
\section*{Question directrice : pourquoi le serveur FS-02 n'était-il pas corrigé~?}
% ============================================================================

Le correctif critique de FS-02 était publié depuis huit mois. Personne ne
l'avait installé~: le serveur ne figurait pas dans l'outil de mise à jour,
et personne n'était chargé de vérifier. De plus, FS-02 avait le même mot
de passe d'administrateur local que onze autres serveurs~: une fois
obtenu, il ouvrait les douze. Ce chapitre traite de la solidité de chaque
machine.

\medskip
\textbf{Plan du chapitre.}
\begin{enumerate}
    \item Réduire la surface d'attaque (section~\ref{sec:c18-surface}).
    \item Durcir Windows et Linux (section~\ref{sec:c18-durcir}).
    \item Gérer les correctifs (section~\ref{sec:c18-correctifs}).
    \item Journaliser (section~\ref{sec:c18-journaux}).
    \item Antivirus, EDR et liste blanche (section~\ref{sec:c18-edr}).
\end{enumerate}

\separateur

% ============================================================================
\section{Réduire la Surface d'Attaque}
\label{sec:c18-surface}
% ============================================================================

\begin{boxdef}
\termecle{Surface d'attaque}~: ensemble des points par lesquels un
attaquant peut interagir avec un système~: services réseau, comptes,
logiciels installés, interfaces. Le \termecle{durcissement} la réduit en
supprimant ce qui est inutile et en configurant sûrement ce qui reste.
\end{boxdef}

Le NIST recommande de partir d'une installation minimale, de retirer les
services, comptes et logiciels inutiles, et de documenter la
configuration retenue~\cite{nist800123}. Les CIS Benchmarks donnent, pour
chaque système, des configurations de référence détaillées et
vérifiables~\cite{cis2021controls}.

% ============================================================================
\section{Durcir Windows et Linux}
\label{sec:c18-durcir}
% ============================================================================

\begin{tableaucours}{L{4.2cm}L{8.0cm}}{%
    \tableauHeader{Mesure} &
    \tableauHeader{Effet}
}
Mot de passe d'administrateur local unique par machine (outil de type
LAPS) & Empêche qu'un mot de passe obtenu sur une machine ouvre toutes les
autres \\
Désactivation de SMBv1 et des protocoles anciens & Ferme les vecteurs de
WannaCry et NotPetya (Chapitre~\ref{chap:attaques}) \\
Stratégies de groupe (GPO) & Applique une configuration commune à toutes
les machines du domaine \\
Pas de droits d'administrateur pour les utilisateurs & Limite l'installation
de logiciels malveillants \\
Pare-feu local actif & Filtre même à l'intérieur d'une zone \\
\end{tableaucours}
\vspace{-0.8em}
\begin{center}{\small\itshape Tableau~18.1. Durcir Windows (exemples)}\end{center}

\begin{boxoutils}{Vérifier quelques points sur un serveur Linux}
\begin{lstlisting}
# connexion SSH : pas de root, pas de mot de passe (clés seulement)
grep -Ei '^(PermitRootLogin|PasswordAuthentication)' /etc/ssh/sshd_config
# qui peut devenir root, et sans mot de passe ?
grep -v '^#' /etc/sudoers /etc/sudoers.d/* | grep -E 'ALL|NOPASSWD'
# fichiers modifiables par tous dans /etc
find /etc -type f -perm -o+w
# services à l'écoute
ss -tlnp
\end{lstlisting}
Chaque ligne affichée appelle une question~: ce service est-il utile~?
cette règle \texttt{NOPASSWD} est-elle justifiée~?
\end{boxoutils}

\begin{boxalert}
\textbf{Violation typique.} Durcir une machine à son installation, puis
ne plus jamais vérifier.

\smallskip
\textbf{Règle~:} une configuration dérive (un administrateur pressé, une
mise à jour, un logiciel ajouté). On la contrôle régulièrement, de
préférence automatiquement, contre la référence.
\end{boxalert}

% ============================================================================
\section{Gérer les Correctifs}
\label{sec:c18-correctifs}
% ============================================================================

Le NIST décrit la gestion des correctifs comme une maintenance préventive,
à planifier selon le risque~\cite{nist80040}.

\begin{boxprocedure}
\textbf{Le processus de gestion des correctifs}
\begin{enumerate}
    \item \textbf{Inventorier} les systèmes et leurs logiciels, avec leurs
          versions (Chapitre~\ref{chap:actifs}).
    \item \textbf{Veiller}~: avis des éditeurs, CVE, alertes nationales.
    \item \textbf{Évaluer}~: gravité (CVSS), exposition, valeur de l'actif,
          exploitation connue.
    \item \textbf{Tester} sur un échantillon représentatif.
    \item \textbf{Déployer} dans le délai fixé pour la catégorie.
    \item \textbf{Vérifier} que le correctif est installé partout.
    \item \textbf{Compenser} quand on ne peut pas corriger (isolement,
          filtrage, surveillance renforcée), sur décision écrite.
\end{enumerate}
\end{boxprocedure}

\begin{tableaucours}{L{5.4cm}L{3.0cm}L{3.8cm}}{%
    \tableauHeader{Catégorie} &
    \tableauHeader{Délai cible} &
    \tableauHeader{Exemple}
}
Critique, exposée à Internet, exploitée & 72~heures & Passerelle d'accès
distant \\
Critique interne & 15~jours & Serveur de fichiers FS-02 \\
Importante & 30~jours & Postes de travail \\
Moyenne et faible & Cycle normal & Logiciels bureautiques \\
\end{tableaucours}
\vspace{-0.8em}
\begin{center}{\small\itshape Tableau~18.2. Délais de correction proposés pour
LiZbiZ (convention à valider par la direction)}\end{center}

% ============================================================================
\section{Journaliser}
\label{sec:c18-journaux}
% ============================================================================

Sans journaux, on ne détecte pas et on ne prouve pas. Le NIST SP~800-92
décrit la gestion des journaux~: génération, transmission, stockage,
analyse et conservation~\cite{nist80092}.

\begin{tableaucours}{L{4.4cm}L{7.8cm}}{%
    \tableauHeader{Événement} &
    \tableauHeader{Pourquoi}
}
Connexions réussies et échouées & Détecter la force brute et l'usage
anormal d'un compte \\
Création et modification de comptes, de groupes & Détecter l'élévation
de privilèges \\
Accès distants & Surveiller la porte d'entrée de l'Opération GAROUA \\
Lancement de processus & Repérer les outils d'attaque \\
Modification des règles de sécurité & Détecter la désactivation des
défenses \\
Accès aux données sensibles & Répondre à « qui a lu quoi~? » \\
\end{tableaucours}
\vspace{-0.8em}
\begin{center}{\small\itshape Tableau~18.3. Les événements à journaliser}\end{center}

Trois conditions rendent les journaux utilisables~: des horloges
synchronisées (NTP), une copie hors de la machine (un attaquant
administrateur efface les journaux locaux), et une durée de conservation
conforme à la loi (Chapitre~\ref{chap:droit}).

% ============================================================================
\section{Antivirus, EDR et Liste Blanche}
\label{sec:c18-edr}
% ============================================================================

\begin{tableaucours}{L{3.2cm}L{4.6cm}L{4.4cm}}{%
    \tableauHeader{Outil} &
    \tableauHeader{Principe} &
    \tableauHeader{Limite}
}
Antivirus & Signatures de fichiers malveillants connus & Aveugle aux outils
légitimes détournés \\
EDR & Enregistre le comportement des machines, détecte et permet de
répondre (isoler, tuer un processus) & Exige une équipe qui traite les
alertes \\
Liste blanche & N'autorise que les programmes approuvés & Lourde à
maintenir sur des postes variés \\
\end{tableaucours}
\vspace{-0.8em}
\begin{center}{\small\itshape Tableau~18.4. Protéger les postes et les serveurs}\end{center}

\begin{boxscenario}{Opération GAROUA. Douze serveurs, un mot de passe}
\textbf{Constatation.} FS-02 et onze autres serveurs partageaient le même
mot de passe d'administrateur local, défini lors de leur installation en
2023 à partir d'une même image.

\medskip
\textbf{Analyse} (confiance élevée). La compromission d'un seul serveur a
donné l'accès aux douze. Un mot de passe local unique et aléatoire par
machine, géré automatiquement, aurait confiné l'attaquant à FS-02.
\end{boxscenario}

\begin{boxmetier}{Administrateur systèmes et réseaux}
\textbf{Mesurer sa propre hygiène.} L'administrateur publie chaque mois
le taux de systèmes conformes à la référence, le taux de correctifs
critiques installés dans les délais, et la liste des exceptions
compensées.
\end{boxmetier}

\begin{boxmetier}{Analyste SOC}
\textbf{Exploiter l'EDR.} L'analyste traite les alertes de l'EDR, isole
une machine suspecte à distance, et recherche sur l'ensemble du parc les
mêmes indicateurs.
\end{boxmetier}

\separateur

\begin{boxresume}
\textbf{Ce qu'il faut retenir de ce chapitre~:}
\begin{itemize}
    \item \textbf{Le durcissement} réduit la surface d'attaque~; les CIS
          Benchmarks donnent des références vérifiables.
    \item \textbf{Un mot de passe local unique par machine} empêche la
          propagation d'une compromission.
    \item \textbf{Les correctifs} suivent un processus et des délais
          fondés sur le risque~; ce qui ne peut être corrigé est compensé.
    \item \textbf{Les journaux} utiles sont choisis, horodatés, copiés hors
          de la machine et conservés.
    \item \textbf{Antivirus, EDR et liste blanche} se complètent~; l'EDR
          exige une équipe.
\end{itemize}
\end{boxresume}

\begin{boxinfo}
\textbf{Points de contrôle de l'élève.} Je sais~: définir la surface
d'attaque~; vérifier la configuration SSH et sudo d'un serveur Linux~;
dérouler le processus de gestion des correctifs~; fixer des délais par
catégorie~; dire quels événements journaliser et pourquoi les copier hors
de la machine.
\end{boxinfo}

\bigskip

\begin{boxexo}[title={\ding{45}\quad Exercice~18.1. Surface d'attaque \niveaubase}]
Listez cinq éléments de la surface d'attaque d'un serveur web, et une
mesure de réduction pour chacun.
\end{boxexo}

\begin{boxexo}[title={\ding{45}\quad Exercice~18.2. SSH \niveaubase}]
Que signifient \texttt{PermitRootLogin no} et
\texttt{PasswordAuthentication no}~? Pourquoi les recommander~?
\end{boxexo}

\begin{boxexo}[title={\ding{45}\quad Exercice~18.3. Les délais \niveaubase}]
Classez dans les catégories du Tableau~18.2~: une faille critique du
portail de suivi, exploitée sur Internet~; une faille importante du
navigateur des postes~; une faille critique d'un serveur interne de test.
\end{boxexo}

\begin{boxexo}[title={\ding{45}\quad Exercice~18.4. sudoers \niveaumoyen}]
La ligne \texttt{technicien ALL=(ALL) NOPASSWD: ALL} figure dans le
fichier \texttt{sudoers} d'un serveur. Analysez le risque et proposez une
règle plus restrictive.
\end{boxexo}

\begin{boxexo}[title={\ding{45}\quad Exercice~18.5. Compenser \niveaumoyen}]
Le logiciel de pesage d'un entrepôt ne fonctionne que sous une version de
Windows qui ne reçoit plus de correctifs. Proposez des mesures
compensatoires.
\end{boxexo}

\begin{boxexo}[title={\ding{45}\quad Exercice~18.6. Les journaux \niveaumoyen}]
Pour l'Opération GAROUA, indiquez quels journaux auraient permis de voir
chacune des étapes, et où ils auraient dû être copiés.
\end{boxexo}

\begin{boxexo}[title={\ding{45}\quad Exercice~18.7. Horloges \niveaumoyen}]
Les serveurs de Garoua avaient sept minutes de retard sur ceux de Douala.
Quelles conséquences pour l'enquête~? Quelle mesure~?
\end{boxexo}

\begin{boxexo}[title={\ding{45}\quad Exercice~18.8. Mesurer \niveaumoyen}]
Rédigez la fiche de l'indicateur « taux de correctifs critiques installés
dans les délais »~: formule, source, fréquence, cible.
\end{boxexo}

\begin{boxexo}[title={\ding{45}\quad Exercice~18.9. Un standard de durcissement \niveauexpert}]
Rédigez le standard de durcissement des serveurs Linux de LiZbiZ, en douze
règles vérifiables, chacune avec sa commande de contrôle.
\end{boxexo}

\begin{boxexo}[title={\ding{45}\quad Exercice~18.10. La politique de correctifs \niveauexpert}]
Rédigez la politique de gestion des correctifs de la BANQUE SANAGA~: rôles,
processus, délais, exceptions, indicateurs.
\end{boxexo}

\bigskip

\begin{boxinfo}
\textbf{Pour aller plus loin.}
\begin{itemize}
    \item NIST SP~800-40 (correctifs), SP~800-92 (journaux), SP~800-123
          (serveurs)~\cite{nist80040,nist80092,nist800123}.
    \item CIS Controls et CIS Benchmarks~\cite{cis2021controls}.
\end{itemize}
\end{boxinfo}

\begin{boxinfo}
\textbf{Transition.} Un système durci peut encore exécuter une
application vulnérable. Le Chapitre~\ref{chap:applications} traite de la
sécurité des applications.
\end{boxinfo}
